\documentclass[11pt]{article}

\usepackage[T1]{fontenc}
\usepackage[utf8]{inputenc}
\usepackage[margin=1.02in]{geometry}
\usepackage{amsmath,amssymb,amsthm,mathtools}
\usepackage{array}
\usepackage{enumitem}
\usepackage{microtype}
\usepackage[hidelinks]{hyperref}

\newtheorem{theorem}{Theorem}
\newtheorem{lemma}{Lemma}
\newtheorem{proposition}{Proposition}
\newtheorem{definition}{Definition}
\newtheorem{claim}{Claim}
\newtheorem{remark}{Remark}
\newtheorem{corollary}{Corollary}

\newcommand{\NS}{\mathrm{NS}}
\newcommand{\Pack}{\mathrm{Pack}}
\newcommand{\Part}{\mathrm{Part}}
\newcommand{\Field}{\mathrm{Field}}
\newcommand{\Member}{\mathrm{Member}}
\newcommand{\Exit}{\mathrm{Exit}}
\newcommand{\CM}{\mathrm{CM}}
\newcommand{\Legal}{\mathrm{Legal}}
\newcommand{\Owork}{\mathfrak O_{\NS}^{\mathrm{work}}}
\newcommand{\WitnessRecord}{\mathcal W}
\newcommand{\Rterm}{\mathcal R_*}
\newcommand{\Ffail}{\mathcal F}
\newcommand{\Pass}{\mathrm{Pass}}
\newcommand{\Fail}{\mathrm{Fail}}
\newcolumntype{L}[1]{>{\arraybackslash}p{#1}}
\setlength{\emergencystretch}{1em}

\title{A Class-Membership Contrapositive Route for the Navier--Stokes Millennium Problem}
\author{Thomas Birnie}
\date{Working manuscript\\May 25, 2026}

\begin{document}
\maketitle

\begin{abstract}
The usual route to the Navier--Stokes Millennium problem is forward-positive.
It tries to prove enough a priori control to continue the smooth solution forever.
That route remains the gold standard because it proves smoothness directly.
It also produces a long list of hard terminal obstructions: source concentration, zero-radius residues, retained-amplitude growth, height flux, signed current, square reserve, endpoint readout, no-pulse failure, and comparison branches.

The route developed here is class-membership contrapositive.
The same smooth-data Navier--Stokes evolution is organized as a working class object.
A legal continuation branch is a member branch.
A finite terminal branch capable of nonsmoothness is tested as a same-solution terminal witness.
The witness must retain three services: a positive same-fluid carrier, the same pressure-viscosity participation tower, and positive-scale one-field coherence.
These are denoted by \(\Pack_Q\), \(\Part_{N,Q}\), and \(\Field_{N,r,Q}\).
The proof shows that every same-solution terminal witness loses one of these services in first-face order.
The lost service is class exit:
\[
  \Exit(Q;\Owork):=\neg\Member(Q;\Owork).
\]
The pass branch remains a member branch and reads out to the usual \(H^s\), \(s>5/2\), continuation theorem.
The fail branch exits the class of legal same-solution continuation candidates.
This is the silver-standard route: the pass-side control mechanism may vary, while every breakdown-capable obstruction leaves the class object whose members are the legal Clay continuation branches.
\end{abstract}

\setcounter{tocdepth}{1}
\tableofcontents
\clearpage

\section{The Problem And The Change Of Question}

The incompressible three-dimensional Navier--Stokes equations are
\[
  \partial_t u +(u\cdot\nabla)u+\nabla p=\nu\Delta u,\qquad
  \nabla\cdot u=0,
\]
with smooth divergence-free initial data \(u_0\).
The Clay problem asks whether the corresponding smooth solution exists for all time.

The direct proof question is:
\[
  \text{which estimate forces the smooth solution to continue?}
\]
That question is natural, and it is the forward-positive gold standard.
One proves a bound strong enough for classical continuation, usually through a continuation norm such as \(H^s\), \(s>5/2\).

The route here changes the question after the forward proof reaches a genuine obstruction.
It asks:
\[
  \text{can a branch capable of nonsmooth terminal failure remain a legal branch of the same class object?}
\]
The two questions have different burdens.
The forward question tries to erase the bad object.
The class-membership question lets the bad object appear and asks whether it can still be the same kind of object the Clay problem asks us to continue.

The basic split is this.
Let \(O\) be a same-solution obstruction reached by a forward Navier--Stokes proof.
The obstruction has a pass side and a fail side:
\[
  O_{\Pass}\Rightarrow \text{the same continuation services remain},
\]
\[
  O_{\Fail}\Rightarrow \text{the branch is capable of finite terminal nonsmoothness}.
\]
The pass side is the member branch.
The fail side is tested as a terminal witness and must lose one of the required services.
Once it loses a required service, it exits the class:
\[
  O_{\Fail}\Rightarrow \Exit(Q;\Owork).
\]
Thus the proof can close without naming in advance which positive estimate saves the pass side.
The fail side must leave the legal same-solution class.

\section{The Forward Obstructions}

The forward-positive program remains the right way to find the hard objects.
It asks each obstruction to be controlled directly.
The present proof uses those obstructions as terminal test objects.

\subsection{Source Control}

Source control tries to prevent terminal concentration of the nonlinear source tree.
Its strongest form is a scale-critical Carleson-type control on donor refill, source reserve, and active ancestry.
The obstruction is that first-moment or local-energy control can leave a square or scale-critical reserve unpaid.

In the class-membership route, selected unpaid donor ancestry is a carrier question first.
A branch whose terminal ancestry has no positive same-fluid carrier fails \(\Pack_Q\).
A retained source residue with a carrier then faces the participation question.
Only after carrier and participation survive may a retained source residue become a Field-readout question.

\subsection{Zero-Radius And Zeno Objects}

Zeno language describes terminal atoms, zero heat-time residues, temporal non-atomicity, and rigid residue classes.
A forward proof tries to rule out those objects positively.
The class-membership proof asks where the selected terminal object lives.

A zero-radius terminal object has no positive terminal carrier.
It fails \(\Pack_Q\).
A retained positive-scale Zeno object stays inside the same Pack/Part/Field witness tree and is read by first failed service.
Zeno is a presentation of terminal behaviour inside the three-service tree.

\subsection{Retained Amplitude And Low Strain}

Retained-amplitude arguments try to turn surviving carrier and participation structure into enough amplitude control for endpoint coherence.
The obstruction is that bounded deformation data can coexist with accumulated positive low-strain action.
Height-flux and low-strain routes sharpen the same pressure point: near-band source charge and frame variation return to the signed source balance.

The class-membership route reads this family as a finite routing engine.
The selected branch lands as Pack gain failure, Part participation or dwell failure, legal retained Field failure, or a Zeno subcase already sorted by the witness tree.

\subsection{Signed Current}

Signed-current arguments try to rectify terminal source-current motion into a no-free-sink or parabolic edge-resistance law.
The obstruction is the variable sign of the strain-production term and the scale-critical flux-amplitude remainder.
Those terms prevent a free coercion theorem.

For class membership, signed current has no independent face.
It returns to Pack, Part, retained Field after licensing, Zeno, or positive supplier support.
Its value is that it identifies where the same services fail.

\subsection{Square Reserve}

Square-reserve arguments ask for a square-source charge law strong enough to control terminal reserve.
They matter because a first-moment source tail can vanish while a square reserve remains active.
The forward branch seeks a direct evolution or charge theorem.

The class-membership branch asks what the selected reserve failure is.
Detached reserve ancestry is Pack-side.
A retained reserve with carrier and participation can become a Field diagnostic.
An absent square-charge identity, by itself, is supplier work.

\subsection{Endpoint Readout And No-Pulse}

Endpoint readout asks for a good scale, a read cover, a coherent field, and an \(H^s\) continuation bound.
No-pulse and averaged routes try to prevent terminal pulses, averaged jumps, and terminal-tail escape.
These are close to classical continuation, so they can sound like the whole proof.

They are downstream in the class-membership argument.
Readout explains why a surviving Pack/Part/Field packet continues.
No-pulse can support retained Field or endpoint production.
Neither replaces the primitive class test.

\subsection{Local Critical Translators And Comparison Branches}

Local \(L^3\) and Duhamel translators are useful when a public-critical-looking terminal mass must be brought back to the same witness record.
They turn a local packet into native source work or a Pack/Part/Field exit.
Their job is branch-local.

Fixed-viscosity Euler-to-Navier--Stokes transfer and Euler mirror comparisons serve different theorem programs.
They illuminate where viscosity enters and where inviscid comparison diverges.
They enter this proof only through a theorem that lands the exact Navier--Stokes object in the same Pack/Part/Field class test.

\section{The Working Class Object}

The proof needs a class object because the contrapositive is a class-law argument.
The class is built from the same datum, same equation, same pressure, same viscosity, and same fluid branch.

\begin{definition}[Working Clay class object]
The working Clay class object \(\Owork\) is the class of same-solution Navier--Stokes continuation packets generated by the original smooth datum and read on one terminal witness record.
For a terminal packet \(Q\), the predicate
\[
  \Member(Q;\Owork)
\]
means that \(Q\) retains the services needed to be a legal same-solution continuation candidate.
The class-exit predicate is
\[
  \Exit(Q;\Owork):=\neg\Member(Q;\Owork).
\]
\end{definition}

\begin{definition}[Witness record]
A witness record is the common record
\[
  \WitnessRecord=(u_0,\nu,u,p,Q,N,r,\Phi,T_*),
\]
where \(u_0\) is the original datum, \((u,p)\) is the same Navier--Stokes solution on \([0,T_*)\), \(Q\) is the terminal packet, \(N\) is the participation depth, \(r\) is the positive field scale, \(\Phi\) is the same-fluid carrier family, and \(T_*\) is the alleged terminal time.
\end{definition}

The common-record condition is important.
A carrier, participation tower, and readout synthesize into one continuation packet only when they belong to the same terminal object.
The services must belong to the same terminal object.

\begin{definition}[The three services]
On a witness record \(\WitnessRecord\), the services are:
\[
  \Pack_Q,\qquad \Part_{N,Q},\qquad \Field_{N,r,Q}.
\]
\(\Pack_Q\) is the positive same-fluid carrier service.
It says that the selected terminal object has a usable positive carrier, cover, or packing geometry for the same fluid.

\(\Part_{N,Q}\) is the same-equation participation service.
It says that on that carrier the differentiated pressure-viscosity tower remains the same Navier--Stokes tower through depth \(N\):
\[
  D_tU_k=K_k+B_k,\qquad
  K_k=-\nabla^{k+1}p+\nu\Delta U_k,\qquad 0\le k\le N.
\]

\(\Field_{N,r,Q}\) is the positive-scale one-field coherence service.
It says that after carrier and participation survive, the terminal packet has coherent field data on a positive scale \(r\) suitable for endpoint readout.
\end{definition}

\begin{definition}[CM packet]
The class-membership packet at depth \(N\), scale \(r\), and terminal packet \(Q\) is
\[
  \CM_{N,r,Q}
  :=
  \Pack_Q\wedge\Part_{N,Q}\wedge\Field_{N,r,Q}
\]
on one witness record.
\end{definition}

\begin{lemma}[Membership readout]
Let \(N_s\) be a depth high enough for the standard \(H^s\), \(s>5/2\), continuation criterion.
If a legal same-solution terminal packet remains a member of \(\Owork\), then
\[
  \Member(Q;\Owork)
  \Longrightarrow
  \sup_{t<T_*}\|u(t)\|_{H^s}<\infty
\]
for some \(s>5/2\), and the same Navier--Stokes solution extends classically beyond \(T_*\).
\end{lemma}

\begin{proof}
Membership means that the terminal packet retains the services needed for a continuation readout on the same witness record.
\(\Pack_Q\) supplies the positive same-fluid carrier.
\(\Part_{N_s,Q}\) supplies the same pressure-viscosity tower through the continuation depth.
\(\Field_{N_s,r,Q}\) supplies positive-scale coherent field data for endpoint readout.
Together these give the displayed \(H^s\) bound at a supercritical continuation level \(s>5/2\).
Classical local well-posedness then continues the same solution past \(T_*\).
\end{proof}

\section{Which Objects Are Tested}

The class-membership proof does not test arbitrary rough objects.
It tests a finite terminal object offered as the Clay breakdown witness for the same original Navier--Stokes solution.

\begin{definition}[CM-test-admissible Clay terminal object]
A finite terminal object \(W_*\) is CM-test-admissible when it is offered as a breakdown witness for the same original smooth-data Navier--Stokes solution.
This means:
\begin{enumerate}[leftmargin=2em]
\item \(W_*\) is attached to the maximal classical solution generated by the original datum \(u_0\).
\item On every preterminal window, the object solves the original incompressible Navier--Stokes system with the same velocity and pressure.
\item The terminal extraction uses the same transported fluid family supplied by the preterminal smooth windows.
\item After legal boundary losses, cutoff losses, non-selected branches, readout artifacts, and paid finite terms are removed, \(W_*\) is still the object claiming to be the finite terminal obstruction.
\item The remaining terminal services available to the object are exactly \(\Pack_Q\), \(\Part_{N,Q}\), and \(\Field_{N,r,Q}\).
\end{enumerate}
\end{definition}

\begin{remark}
CM-test-admissible means admissible for the test.
It does not mean member.
Membership or exit is the conclusion.
\end{remark}

\begin{lemma}[Finite Clay breakdown enters the test]
Every finite same-surface Clay breakdown witness is CM-test-admissible.
\end{lemma}

\begin{proof}
A finite same-surface Clay breakdown witness is, by definition, the original smooth-data Navier--Stokes solution stopped at a finite terminal time where classical continuation is alleged to fail.
The datum, equation, pressure, velocity, viscosity, and preterminal fluid family are therefore the same ones used before the terminal time.
Legal losses and irrelevant branches can be stripped away without changing the alleged obstruction.
The remaining object is precisely the terminal object to be tested by the three services.
\end{proof}

\section{Terminal Packet Capture}

The capture theorem is the bridge from a hypothetical finite breakdown to the Pack-first decision tree.

\begin{theorem}[Terminal packet capture]
Let \((u,p)\) be the maximal smooth Navier--Stokes solution generated by \(u_0\) on \([0,T_*)\), with \(T_*<\infty\) assumed.
Let \(\Rterm\) be a same-solution terminal obstruction after legal losses, non-selected branches, readout artifacts, and paid finite terms have been removed.
Fix \(N_s\) high enough for \(H^s\), \(s>5/2\), continuation readout.
Then
\[
  \Rterm
  \Longrightarrow
  \neg\Pack_Q
  \vee
  \neg\Part_{N_s,Q}
  \vee
  \forall r>0\,\neg\Field_{N_s,r,Q}.
\]
Equivalently,
\[
  \Rterm
  \Longrightarrow
  \begin{cases}
  \neg\Pack_Q,\\
  \Pack_Q\wedge\neg\Part_{N_s,Q},\\
  \Pack_Q\wedge\Part_{N_s,Q}\wedge\forall r>0\,\neg\Field_{N_s,r,Q}.
  \end{cases}
\]
\end{theorem}

\begin{proof}
For every \(\varepsilon>0\), the restriction of \((u,p)\) to \([0,T_*-\varepsilon]\) is classical.
Each compact preterminal window carries the same datum, same equation, same pressure, same viscosity, same transported fluid family, and hence the same local continuation services.

Pass to a terminal subsequence producing \(\Rterm\).
The selected object first asks whether a positive same-fluid carrier survives.
If no such carrier survives, \(\neg\Pack_Q\).

Assume that \(\Pack_Q\) survives.
The next question is whether the same differentiated Navier--Stokes pressure-viscosity tower survives on that carrier through depth \(N_s\).
If the tower does not survive, \(\neg\Part_{N_s,Q}\).

Assume that both \(\Pack_Q\) and \(\Part_{N_s,Q}\) survive.
If there is a positive scale \(r>0\) with \(\Field_{N_s,r,Q}\), then the full packet \(\CM_{N_s,r,Q}\) survives on the same ledger.
By the membership readout lemma, the same solution has a bounded \(H^s\), \(s>5/2\), continuation norm and extends classically past \(T_*\).
That contradicts the terminal role of \(\Rterm\).
Thus \(\Field_{N_s,r,Q}\) fails for every \(r>0\).
\end{proof}

\begin{corollary}[No fourth primitive service]
After legal losses, readout artifacts, non-selected branches, and paid finite terms are removed, a same-solution terminal obstruction has no primitive service outside Pack, Part, and Field.
\end{corollary}

\begin{proof}
The preceding theorem tests exactly the services needed for the same-solution continuation packet.
Any proposed residue either lacks the carrier, lacks participation after carrier survives, lacks positive-scale coherence after carrier and participation survive, or is a supplier/readout description waiting for one of those services to be tested.
\end{proof}

\section{Branch-Family Exhaustion}

The previous theorem explains the proof spine.
This section records how the current branch field is represented without turning the paper into an inventory.
Each family enters through the same question: what terminal object is selected, and which service fails first?

\begin{proposition}[Branch-family exhaustion]
Every current non-Euler Navier--Stokes branch family in the working class-membership program has one of the following mathematical roles:
\[
  \Pack\text{-exit},\quad
  \Part\text{-exit},\quad
  \Field\text{-exit},\quad
  \text{finite typed disjunction},\quad
  \text{support until promoted by a named face theorem}.
\]
\end{proposition}

\begin{proof}
We read the families by their selected terminal object.

Source-control and donor-refill branches select terminal source ancestry or reserve.
Detached or unpaid ancestry without a positive same-fluid carrier is Pack-exit.
Retained source residue with carrier and broken pressure-viscosity participation is Part-exit.
Retained source residue that has passed Pack and Part becomes a Field/readout diagnostic.

Zero-radius and terminal-atom branches select an object with no positive terminal carrier unless a retained positive-scale packet has first been supplied.
The zero-radius branch is Pack-exit.
Retained Zeno alternatives are finite typed subcases within the same witness tree.

Retained-amplitude, height-flux, low-strain, signed-current, and square-reserve branches select failures of gain, dwell, participation, source balance, or retained coherence.
Each such failure either breaks the carrier, breaks the same pressure-viscosity participation tower after the carrier survives, breaks positive-scale coherence after Pack and Part survive, or remains a positive supplier theorem outside the class primitive.

Good-scale, endpoint-readout, no-pulse, averaged, and terminal-tail branches read a packet that has already retained enough structure to be read.
They support the Field and continuation consequences.
They become proof material only when the exact packet has already entered Pack, Part, or Field.

Local critical translators place local \(L^3\) or Duhamel mass back on the same witness ledger.
Their role is branch-local: native source work or Pack/Part/Field exit.

Fixed-viscosity transfer and Euler mirror branches are comparison programs.
They remain outside the non-Euler class-membership proof until a theorem transfers the exact Navier--Stokes object into the same witness tree.

These alternatives exhaust the branch families because every branch either supplies a selected same-solution terminal object or supplies context for such an object.
Selected terminal objects are exhausted by terminal packet capture.
Context without a selected face remains support.
\end{proof}

\begin{center}
\small
\begin{tabular}{L{0.30\linewidth}L{0.31\linewidth}L{0.30\linewidth}}
\textbf{Branch family} & \textbf{Forward-positive burden} & \textbf{CM role}\\
\hline
Source control and donor refill & prove scale-critical source reserve control & Pack first; Part after carrier survives; retained Field only after Pack and Part\\
Zero-radius and Zeno & prove temporal non-atomicity or rigid residue exclusion & zero-radius Pack exit; retained positive-scale typed subcases\\
Retained amplitude and low strain & prove gain, BV, or height-flux control & Pack/Part/Field/Zeno finite routing\\
Signed current & prove no-free-sink or edge resistance & support returning to Pack, Part, retained Field, or Zeno\\
Square reserve & prove square-source charge or native trilinear control & supplier unless selected reserve lands in first-face order\\
Endpoint readout and no-pulse & prove good scale, read cover, and pulse exclusion & downstream support for membership readout and Field\\
Local critical translators & localize public-critical terminal mass & branch-local native source work or class exit\\
Transfer and mirror comparison & compare neighbouring theorem worlds & outside this proof until an exact transfer lands\\
\end{tabular}
\end{center}

\section{Class Exit}

The class-exit theorem turns first-face failure into the contrapositive conclusion.

\begin{theorem}[Finite failure witness exits the class]
Let \(W_*\) be a CM-test-admissible finite Clay terminal object.
Then for a suitable terminal packet \(Q\) and continuation depth \(N_s\),
\[
  W_*
  \Longrightarrow
  \Exit(Q;\Owork).
\]
\end{theorem}

\begin{proof}
By terminal packet capture, \(W_*\) loses a first required service:
\[
  \neg\Pack_Q
  \vee
  \neg\Part_{N_s,Q}
  \vee
  \forall r>0\,\neg\Field_{N_s,r,Q}.
\]
If \(\Pack_Q\) fails, the terminal object has no positive same-fluid carrier and cannot be a member packet.
If \(\Part_{N_s,Q}\) fails after Pack survives, the terminal object no longer participates in the same pressure-viscosity Navier--Stokes tower and cannot be a member packet.
If every positive \(\Field_{N_s,r,Q}\) fails after Pack and Part survive, the terminal object lacks the positive-scale coherence required for membership readout and cannot be a member packet.

Each case is the loss of a required service for \(\Member(Q;\Owork)\).
Therefore
\[
  \Exit(Q;\Owork):=\neg\Member(Q;\Owork).
\]
\end{proof}

\section{The Class-Law Proof}

Classification alone is not the whole argument.
The proof closes at the class law: the legal same-solution continuation candidates are exactly the branches represented inside the working class object.

\begin{proposition}[Legal branch law]
Every legal same-solution continuation branch for the Clay Navier--Stokes problem is represented by the working class object:
\[
  \Legal_{\NS}(Q)\Longrightarrow \Member(Q;\Owork).
\]
\end{proposition}

\begin{proof}
A legal same-solution continuation branch is not an arbitrary terminal distribution.
It is the original smooth-data Navier--Stokes evolution continued as the same velocity-pressure solution with the same viscosity and same fluid participation.
To be such a continuation branch, it must retain a positive same-fluid carrier, the same pressure-viscosity participation tower through the continuation depth, and positive-scale coherent field data for the \(H^s\), \(s>5/2\), readout.
Those are exactly the services defining membership in \(\Owork\).
\end{proof}

\begin{theorem}[CM contrapositive Navier--Stokes class theorem]
Inside the working class formulation, no finite same-surface Clay breakdown branch remains as an in-class counterexample.
More explicitly,
\[
  \text{finite same-surface Clay breakdown}
  \Longrightarrow
  \Exit(Q;\Owork),
\]
while
\[
  \Legal_{\NS}(Q)\Longrightarrow\Member(Q;\Owork)\Longrightarrow\text{smooth continuation}.
\]
Consequently the only legal same-solution branch remaining inside the Clay continuation class is smooth.
\end{theorem}

\begin{proof}
Assume a finite same-surface Clay breakdown branch.
By the finite-breakdown entry lemma, its terminal object is CM-test-admissible.
By terminal packet capture, that object loses a first required service: Pack, Part, or Field.
By finite failure witness class exit, the branch is \(\Exit(Q;\Owork)\).

Now take any branch that remains legal as a same-solution continuation candidate.
By the legal branch law, it is a member of \(\Owork\).
By the membership readout lemma, membership gives an \(H^s\), \(s>5/2\), continuation bound and extends the same solution classically.

Thus the branch capable of finite terminal nonsmoothness leaves the legal same-solution class, and every branch left inside that class continues smoothly.
No finite same-surface Clay breakdown branch survives as an in-class counterexample.
\end{proof}

\section{The Short Proof Spine}

Once the definitions and branch law are in place, the proof is short.

\begin{claim}[Proof spine]
Let \(O\) be any obstruction reached by the forward-positive Navier--Stokes program.
Split it into a pass branch and a fail branch:
\[
  O_{\Pass}:\text{ the same continuation services remain},
\qquad
  O_{\Fail}:\text{ the branch can witness finite nonsmoothness}.
\]
Then
\[
  O_{\Pass}\Rightarrow\Member(Q;\Owork)\Rightarrow\text{smooth continuation},
\]
and
\[
  O_{\Fail}\Rightarrow\Exit(Q;\Owork).
\]
\end{claim}

\begin{proof}
The pass branch retains the same-fluid carrier, same-equation participation, and positive-scale field readout required for \(\Member(Q;\Owork)\).
The membership readout continues the solution.

The fail branch is offered as a finite terminal witness for the same solution.
It is CM-test-admissible and enters the Pack-first terminal tree.
It loses Pack, or loses Part after Pack survives, or loses Field at every positive scale after Pack and Part survive.
The first lost service is class exit.
\end{proof}

\section{Referee Check}

The proof has one load-bearing identification:
\[
  \Legal_{\NS}(Q)\Longrightarrow\Member(Q;\Owork).
\]
That is the class-law hinge.
The classification theorem supplies the exit witness.
The membership readout supplies smooth continuation for legal member branches.
The branch law says why an exit branch is no longer a legal in-class continuation candidate.

A demand to prove one particular source, reserve, no-pulse, or readout estimate belongs to the forward-positive program.
Such a theorem can strengthen the pass side and simplify the branch field.
The class-membership contrapositive needs a different statement: every failure branch capable of finite nonsmoothness must be canonical enough to be tested, must lose a required service, and must exit the class.

This is the paper's mathematical promise.
The forward program shows where the real obstructions arise.
The class object supplies the legal branch test.
The terminal packet capture exhausts the fail side.
The membership readout continues the pass side.

\end{document}
